% paper_prl_proton_tagged.tex -- Letter: the proton-tagged measurement.
% Companion to paper_prd_inclusive.tex (joint submission; see the cover-letter note at the end).
% PRL limit 3750 words including figure word-equivalents; see report/tools/prl_wordcount.py.
\documentclass[aps,prl,twocolumn,superscriptaddress,nofootinbib,floatfix]{revtex4-2}
\usepackage{amsmath,amssymb}
\usepackage{graphicx}
\usepackage{booktabs}
\usepackage{xcolor}
\usepackage{hyperref}
\graphicspath{{figures/}{figures/bkgfit/}}
% Shared macros for the two papers. Kept identical to the analysis note's definitions.
\newcommand{\numu}{$\nu_\mu$}
\newcommand{\numubar}{$\bar{\nu}_\mu$}
\newcommand{\ccpip}{$\mathrm{CC}\pi^{\pm}$}
\newcommand{\uboone}{MicroBooNE}
\newcommand{\GeVc}{\,\text{GeV}\!/c}
\newcommand{\pot}{\,\text{POT}}
\newcommand{\pmuM}{p_\mu}
\newcommand{\ppiM}{p_\pi}
\newcommand{\cthmuM}{\cos\theta_\mu}
\newcommand{\cthpiM}{\cos\theta_\pi}
\newcommand{\thmupM}{\theta_{\mu\pi}}
\newcommand{\AC}{A_C}
% Draft status. Every data point in both papers is a Poisson throw of the central-value
% simulation: the signal region is blind. FHC and combined normalisations are being rebuilt on
% the corrected Run-1 exposure and the numbers here predate that rebuild.
\newcommand{\draftbanner}{\begin{center}\fbox{\parbox{0.95\linewidth}{\small\textbf{DRAFT ---
  blind analysis.} Every ``data'' point is a Poisson throw of the central-value simulation. Numbers
  are placeholders for the unblinded result; FHC and combined normalisations predate the Run-1
  exposure correction and will move by $\approx\!+14\%$ and $+6\%$.}}\end{center}}


\begin{document}

\title{First measurement of transverse kinematic imbalance and hadronic-system kinematics in
  charged-current single-pion production on argon}

\author{The \uboone{} Collaboration}
\noaffiliation

\date{\today}

\begin{abstract}
We report the first measurement of the kinematics of the full visible hadronic system in
charged-current single-charged-pion production, $\nu_\mu(\bar\nu_\mu)\,\mathrm{Ar}\to
\mu^\mp\pi^\pm p\,X$, using the \uboone{} liquid-argon time-projection chamber in the NuMI beam
with $1.99\times10^{21}$ protons on target in forward- and reverse-horn-current running. Requiring
a reconstructed proton alongside the muon and pion gives access to observables that the inclusive
channel cannot see: the transverse kinematic imbalance between the muon and the $\pi p$ system
($\delta p_T$, $\delta\alpha_T$, $\delta\phi_T$ and the inferred nucleon momentum $p_n$), the
calorimetric hadronic mass $W_{\rm had}$, the proton polar angle and the pion--proton opening
angle. Flux-averaged total and differential cross sections are extracted by Wiener-SVD unfolding
with complete covariance and are compared with GENIE, NuWro, NEUT and GiBUU. The $\delta\alpha_T$
distribution rises toward $180^\circ$, the region most sensitive to final-state interactions, and
the standalone generators span a factor of $1.6$ in the proton-tagged total. [DRAFT: values are
blind placeholders.]
\end{abstract}

\maketitle
\draftbanner

\paragraph*{Introduction.}
Neutrino--nucleus interactions in the few-GeV region are the largest source of systematic
uncertainty in accelerator oscillation experiments, and charged-current single-pion production
(CC$\pi^\pm$) is both a signal channel and the dominant background to quasi-elastic-like samples.
Its cross section on argon is now measured
inclusively~\cite{uboone_bnb2026_ccpi,uboone_nue_ccpi} and is the subject of a companion
article~\cite{companion_prd}. Those measurements characterise the lepton and the pion. They are
blind to the nucleon, and therefore to the quantity that most directly encodes nuclear effects:
the momentum balance between the outgoing lepton and the outgoing hadronic system.

Transverse kinematic imbalance (TKI)~\cite{lu_tki,furmanski_pn} exploits that balance. For a
neutrino incident along $\hat z$ on a free nucleon at rest, the momenta of the lepton and the
hadronic system transverse to the beam cancel; on a nucleus the imbalance $\delta p_T$ measures the
initial-state nucleon momentum, and its orientation $\delta\alpha_T$ relative to the lepton is
flat in the absence of final-state interactions (FSI) and rises toward $180^\circ$ when the
hadrons lose momentum in the nucleus. TKI has been measured in pionless final states by
T2K~\cite{t2k_tki}, MINER$\nu$A~\cite{minerva_tki} and \uboone{}~\cite{uboone_tki}. In a
pion-production channel the hadronic system is $\pi+p$ and the same variables probe FSI acting on
a pion and a proton together, a regime in which generators disagree most and in which no
measurement on argon exists. We present that measurement.

\paragraph*{Experiment and data.}
\uboone{}~\cite{uboone_nim} is an $85$-tonne active-mass liquid-argon time-projection chamber
located $135$~mrad off the axis of the NuMI beam~\cite{numi_beam}. The off-axis flux peaks below
$1$~GeV and, because the detector is off axis, the neutrino direction is $28^\circ$ from the
detector $z$ axis. Every angle in this Letter is defined about the neutrino direction; defining
them about the detector axis would move about $65\%$ of signal events across a $\cos\theta_\mu$
bin~\cite{companion_prd}. The data comprise the full NuMI exposure: $8.86\times10^{20}$~POT in
forward horn current (FHC, $\nu_\mu$-dominated), $1.11\times10^{21}$ in reverse horn current (RHC,
$\bar\nu_\mu$-enriched), and their combination. Each running period is normalised to its own
exposure and each horn configuration carries its own integrated $\nu_\mu+\bar\nu_\mu$ flux, from
the PPFX-corrected simulation~\cite{ppfx}. The simulation uses GENIE v3~\cite{genie3} with the
\uboone{} tune~\cite{uboone_tune}, GEANT4 and the Pandora reconstruction~\cite{pandora,uboone_pandora}.

\paragraph*{Signal and selection.}
The signal is defined by the final state leaving the nucleus: one muon with $p_\mu>0.15\GeVc$,
exactly one charged pion with $p_\pi>0.175\GeVc$, no neutral pion, and at least one proton with
$p_p>0.3\GeVc$, all from an interaction vertex inside the fiducial volume, with muon--pion opening
angle $\theta_{\mu\pi}<2.6$~rad. Both neutrino and antineutrino interactions are signal. The
selection applies the inclusive CC$\pi^\pm$ selection of the companion article~\cite{companion_prd}
--- a muon candidate, exactly one contained pion candidate, an electromagnetic-shower veto and
the opening-angle requirement --- and additionally requires a contained proton candidate,
identified by a log-likelihood-ratio particle-identification score~\cite{uboone_llr}, with
range-based momentum above $0.3\GeVc$. The proton-tagged selection is $11.4\%$ efficient and
$51\%$ pure; the dominant background is CC$0\pi$ with a proton misidentified as the pion,
followed by multi-pion final states. The proton candidate is the leading true proton in $82\%$ of
selected signal events and a true proton in $87\%$; the remainder, where a pion or other track is
taken for the proton, is absorbed into the response and is the dominant source of migration for
the angular observables.

\paragraph*{Observables.}
With $\vec p_{\rm had}=\vec p_\pi+\vec p_p$ and transverse components taken about the neutrino
direction,
\begin{align}
  \delta p_T &= |\vec p_T^{\,\mu}+\vec p_T^{\,\rm had}|, \\
  \delta\alpha_T &= \arccos\frac{-\vec p_T^{\,\mu}\cdot(\vec p_T^{\,\mu}+\vec p_T^{\,\rm had})}
                              {|\vec p_T^{\,\mu}|\,\delta p_T}, \\
  \delta\phi_T &= \arccos\frac{-\vec p_T^{\,\mu}\cdot\vec p_T^{\,\rm had}}
                            {|\vec p_T^{\,\mu}|\,|\vec p_T^{\,\rm had}|},
\end{align}
and $p_n$ is the initial nucleon momentum inferred from $\delta p_T$ and the longitudinal
imbalance under the assumption of a stationary target nucleus~\cite{furmanski_pn}. The
calorimetric hadronic mass is $W_{\rm had}=[m_p^2+2m_p(E_{\rm cal}-E_\mu)-Q^2]^{1/2}$. The proton
polar angle $\theta_p$ is measured about the neutrino direction and the pion--proton opening angle
$\theta_{\pi p}$ between the two reconstructed tracks, so the latter needs no beam direction at
all. The directly reconstructed $\pi p$ invariant mass, which would show the $\Delta(1232)$, is not
measurable as a shape at this exposure: no binning satisfies both the migration and the statistics
requirement, and only its integral is reported.

\paragraph*{Cross-section extraction.}
Differential cross sections are extracted by Wiener-SVD unfolding~\cite{wienervsd} of the
background-subtracted reconstructed spectrum with a response built from the simulation, using the
\textsc{XSecAnalyzer} framework~\cite{xsec_analyzer}. Regularisation leaves a known residual
smearing, described by the additional-smearing matrix $\AC$, which is released with each result: a
prediction must be multiplied by $\AC$ before comparison. Because $\AC$ is not
normalisation-preserving, the integral of a differential result is not the total cross section;
the total is obtained separately from a one-bin extraction in which $\AC\equiv1$. Uncertainties
from the flux (PPFX universes), the interaction model (GENIE reweighting), hadron re-interaction,
detector response (dedicated variation samples~\cite{uboone_detvar}), exposure and target count,
and simulation, beam-off and data statistics are propagated as full covariance matrices through
the same unfolding and released with the results. The proton-tagged sample is statistically
sparse and correspondingly detector-sensitive: the detector term ranges from $10\%$ to over
$100\%$ across bins, with a median of $33\%$, against $16$--$35\%$ inclusively. The binning of every
observable was set by requiring a column-normalised migration diagonal above $0.68$ and a
well-conditioned $\AC$, which is why the four TKI observables are reported in two bins each.

\paragraph*{Results.}
\begin{table}[b]
  \caption{Flux-averaged proton-tagged total cross section from the one-bin extraction,
    $10^{-38}$~cm$^2$/Ar, with the standalone-generator predictions for FHC. [DRAFT: blind
    placeholders; FHC and combined normalisations predate the Run-1 exposure correction.]}
  \label{tab:total1p}
  \begin{ruledtabular}
  \begin{tabular}{lcc}
    Configuration & $\sigma$ & unc.\ [\%] \\
    \hline
    FHC & $0.618\pm0.255$ & $41.3$ \\
    RHC & $0.531\pm0.251$ & $47.2$ \\
    Combined & $0.571\pm0.240$ & $42.1$ \\
    \hline
    GENIE / NuWro / GiBUU / NEUT (FHC) & \multicolumn{2}{c}{$0.50$ / $0.65$ / $0.74$ / $0.78$} \\
  \end{tabular}
  \end{ruledtabular}
\end{table}
Table~\ref{tab:total1p} gives the total cross section in the three configurations. The
standalone generators span a factor of $1.6$ and bracket the result; at $41$--$47\%$ uncertainty
the total does not yet discriminate among them, and the discrimination lies in the shapes.

\begin{figure*}[t]
  \centering
  \includegraphics[width=\textwidth]{dsigma_ccpi1p_noW_COMB}
  \caption{Proton-tagged differential cross sections in the combined FHC$+$RHC configuration:
    the calorimetric hadronic mass $W_{\rm had}$ and the four transverse-kinematic-imbalance
    observables $\delta p_T$, $\delta\alpha_T$, $\delta\phi_T$ and $p_n$. Points are the unfolded
    data with the total uncertainty; the grey histogram is the $\AC$-smeared central-value
    prediction and the coloured histograms are GENIE, GiBUU, NEUT and NuWro, all multiplied by
    $\AC$. [DRAFT: blind fake data.]}
  \label{fig:tki}
\end{figure*}

Figure~\ref{fig:tki} shows the five hadronic observables in the combined configuration.
$\delta\alpha_T$ rises toward $180^\circ$: the enhancement in the backward region is the
signature that the hadronic system has lost transverse momentum inside the nucleus, and it is
sensitive to final-state interactions and other nuclear effects. The four generators differ most
in this bin, in the high-$\delta p_T$ tail and in the upper region of $W_{\rm had}$, where NEUT
and GiBUU lie above the tune and GENIE below it. The differential results close on their own
$\AC$-smeared truth to $0.96$--$1.04$ in every observable with $\chi^2/{\rm ndf}$ below unity.

\begin{figure*}[t]
  \centering
  \includegraphics[width=\textwidth]{newobs_xsec}
  \caption{Proton polar angle $\theta_p$ (top) and pion--proton opening angle $\theta_{\pi p}$
    (bottom) for FHC, RHC and the combined configuration, with the four generator predictions and
    integrated cross sections in the legends. The RHC $\theta_p$ response is rank-deficient and that
    panel is not a shape measurement. [DRAFT: blind fake data.]}
  \label{fig:angles}
\end{figure*}

Figure~\ref{fig:angles} shows the two angular observables. Both distributions are strongly
forward peaked, the top $\theta_p$ bin carrying a fifteenth of the first bin's density, and the
pion and proton are separated by less than $2$~rad in the bulk of events. The generators differ
mainly in normalisation --- GENIE lowest, NEUT highest, a spread of $1.6$ in the integral --- with
NEUT's relative excess in the middle $\theta_{\pi p}$ bin the one shape difference. $\theta_{\pi p}$
is the best-conditioned observable of the proton-tagged set; $\theta_p$ is reported for FHC and
the combined configuration only, its RHC response being rank-deficient at any binning tried.

\paragraph*{Validation.}
The extraction is validated in three ways on simulation~\cite{companion_prd}: closure of
pseudo-data thrown from the nominal model, which tests the implementation; ensembles of
independent Poisson realisations, which return correctly sized statistical uncertainties for the
inclusive extractions; and unfolding of pseudo-data thrown from a materially different interaction
model with the nominal response, in which no inclusive bin is biased by more than $0.8\sigma$ of its
uncertainty. The last two have not yet been run for the proton-tagged sample and are required
before submission. The background model is tested in four control regions disjoint from the
signal region, opened on beam-on data, which indicate no correction to the background-class
normalisations.

\paragraph*{Conclusions.}
We have presented the first measurement of the hadronic-system kinematics in charged-current
single-pion production on argon: the transverse kinematic imbalance, the calorimetric hadronic
mass, the proton polar angle and the pion--proton opening angle, together with the flux-averaged
total, in neutrino-dominated, antineutrino-enriched and combined NuMI exposures. The
$\delta\alpha_T$ enhancement toward $180^\circ$ is observed in a pion-production channel for the
first time, and the standalone generators span a factor of $1.6$ in normalisation with shapes that
differ most where final-state interactions act. The additional-smearing matrices and covariances
for every result are released with this Letter.

\begin{acknowledgments}
[Standard \uboone{} acknowledgments.]
\end{acknowledgments}

\begin{thebibliography}{99}

\bibitem{numi_nue_ccpi}
  MicroBooNE Collaboration,
  ``Measurement of the flux-averaged differential $\nu_e$ CC$1\pi$ cross section
  in the NuMI beam,'' Internal Note v3.2 (2025).

\bibitem{mistry_thesis}
  K.~Mistry,
  ``First Measurement of the Flux-Averaged Differential Charged-Current
  Electron-Neutrino and Antineutrino Cross Section on Argon with the MicroBooNE
  Detector,'' PhD thesis (2021).

\bibitem{nova}
  NOvA Collaboration, M.~A.~Acero \textit{et al.},
  ``New constraints on oscillation parameters from $\nu_e$ appearance and
  $\nu_\mu$ disappearance in the NOvA experiment,''
  \textit{Phys.\ Rev.\ D} \textbf{98}, 032012 (2018).

\bibitem{t2k}
  T2K Collaboration, K.~Abe \textit{et al.},
  ``Constraint on the matter--antimatter symmetry-violating phase in neutrino
  oscillations,''
  \textit{Nature} \textbf{580}, 339--344 (2020).

\bibitem{dune}
  DUNE Collaboration, B.~Abi \textit{et al.},
  ``Deep Underground Neutrino Experiment (DUNE), Far Detector Technical Design
  Report, Volume II: DUNE Physics,'' arXiv:2002.03005 [hep-ex] (2020).

\bibitem{pandora}
  J.~S.~Marshall and M.~A.~Thomson,
  ``The Pandora Software Development Kit for Pattern Recognition,''
  \textit{Eur.\ Phys.\ J.\ C} \textbf{75}, 439 (2015).

\bibitem{wienervsd}
  W.~Tang, X.~Li, X.~Qian, H.~Wei, and C.~Zhang,
  ``Data Unfolding with Wiener-SVD Method,''
  \textit{JINST} \textbf{12}, P10002 (2017), arXiv:1705.03568 [physics.data-an].

\bibitem{ppfx}
  L.~Aliaga \textit{et al.}\ (MINERvA Collaboration),
  ``Neutrino Flux Predictions for the NuMI Beam,''
  \textit{Phys.\ Rev.\ D} \textbf{94}, 092005 (2016).

\bibitem{genie}
  C.~Andreopoulos \textit{et al.},
  ``The GENIE Neutrino Monte Carlo Generator,''
  \textit{Nucl.\ Instrum.\ Methods A} \textbf{614}, 87--104 (2010).

\bibitem{sce}
  MicroBooNE Collaboration, P.~Abratenko \textit{et al.},
  ``Measurement of space charge effects in the MicroBooNE LArTPC using cosmic
  muons,''
  \textit{JINST} \textbf{15}, P12037 (2020).

\bibitem{uboone_laser}
  MicroBooNE Collaboration, C.~Adams \textit{et al.},
  ``A Method to Determine the Electric Field of Liquid Argon Time Projection
  Chambers Using a UV Laser System and Its Application in MicroBooNE,''
  \textit{JINST} \textbf{15}, P07010 (2020).

\bibitem{uboone_pandora}
  MicroBooNE Collaboration, R.~Acciarri \textit{et al.},
  ``The Pandora multi-algorithm approach to automated pattern recognition of
  cosmic-ray muon and neutrino events in the MicroBooNE detector,''
  \textit{Eur.\ Phys.\ J.\ C} \textbf{78}, 82 (2018).

\bibitem{uboone_llr}
  MicroBooNE Collaboration, P.~Abratenko \textit{et al.},
  ``Calorimetric classification of track-like signatures in liquid argon TPCs
  using MicroBooNE data,''
  \textit{JHEP} \textbf{12}, 153 (2021).

\bibitem{uboone_mcs}
  MicroBooNE Collaboration, P.~Abratenko \textit{et al.},
  ``Determination of muon momentum in the MicroBooNE LArTPC using an improved
  model of multiple Coulomb scattering,''
  \textit{JINST} \textbf{12}, P10010 (2017).

\bibitem{uboone_detvar}
  MicroBooNE Collaboration, P.~Abratenko \textit{et al.},
  ``Novel approach for evaluating detector-related uncertainties in a LArTPC
  using MicroBooNE data,''
  \textit{Eur.\ Phys.\ J.\ C} \textbf{82}, 454 (2022).

\bibitem{uboone_tune}
  MicroBooNE Collaboration, P.~Abratenko \textit{et al.},
  ``New CC0$\pi$ GENIE model tune for MicroBooNE,''
  \textit{Phys.\ Rev.\ D} \textbf{105}, 072001 (2022), arXiv:2110.14028 [hep-ex].

\bibitem{uboone_tki}
  MicroBooNE Collaboration, P.~Abratenko \textit{et al.},
  ``First double-differential measurement of kinematic imbalance in neutrino
  interactions with the MicroBooNE detector,''
  \textit{Phys.\ Rev.\ Lett.} \textbf{131}, 101802 (2023), arXiv:2301.03700 [hep-ex].

\bibitem{xsec_analyzer}
  S.~Gardiner \textit{et al.},
  ``XSecAnalyzer: a cross-section analysis framework for MicroBooNE,''
  \url{https://github.com/uboone/xsec_analyzer} (accessed 2026).

\bibitem{numi_beam}
  P.~Adamson \textit{et al.},
  ``The NuMI Neutrino Beam,''
  \textit{Nucl.\ Instrum.\ Methods A} \textbf{806}, 279--306 (2016).

\bibitem{genie3}
  L.~Alvarez-Ruso \textit{et al.}\ (GENIE Collaboration),
  ``Recent highlights from GENIE v3,''
  \textit{Eur.\ Phys.\ J.\ ST} \textbf{230}, 4449--4467 (2021).

\bibitem{nuwro}
  T.~Golan, J.~T.~Sobczyk, and J.~\.Zmuda,
  ``NuWro: the Wroc\l{}aw Monte Carlo generator of neutrino interactions,''
  \textit{Nucl.\ Phys.\ B Proc.\ Suppl.} \textbf{229--232}, 499 (2012).

\bibitem{gibuu}
  O.~Buss \textit{et al.},
  ``Transport-theoretical description of nuclear reactions,''
  \textit{Phys.\ Rept.} \textbf{512}, 1--124 (2012);
  U.~Mosel, ``Neutrino event generators: foundation, status and future,''
  \textit{J.\ Phys.\ G} \textbf{46}, 113001 (2019).

\bibitem{neut2021}
  Y.~Hayato and L.~Pickering,
  ``The NEUT neutrino interaction simulation program library,''
  \textit{Eur.\ Phys.\ J.\ ST} \textbf{230}, 4469--4481 (2021).

\bibitem{geant4}
  S.~Agostinelli \textit{et al.}\ (GEANT4 Collaboration),
  ``GEANT4: a simulation toolkit,''
  \textit{Nucl.\ Instrum.\ Methods A} \textbf{506}, 250--303 (2003).

\bibitem{geant4reweight}
  J.~Calcutt, C.~Thorpe, K.~Mahn, and L.~Fields,
  ``Geant4Reweight: a framework for evaluating and propagating hadronic
  interaction uncertainties in Geant4,''
  \textit{JINST} \textbf{16}, P08042 (2021).

\bibitem{tmva}
  A.~Hoecker \textit{et al.},
  ``TMVA: Toolkit for Multivariate Data Analysis,''
  PoS \textbf{ACAT}, 040 (2007), arXiv:physics/0703039.

\bibitem{dagostini}
  G.~D'Agostini,
  ``A multidimensional unfolding method based on Bayes' theorem,''
  \textit{Nucl.\ Instrum.\ Methods A} \textbf{362}, 487--498 (1995).

\bibitem{lu_tki}
  X.-G.~Lu, L.~Pickering, S.~Dolan, G.~Barr, D.~Coplowe, Y.~Uchida, D.~Wark,
  M.~O.~Wascko, A.~Weber, and T.~Yuan,
  ``Measurement of nuclear effects in neutrino interactions with minimal
  dependence on neutrino energy,''
  \textit{Phys.\ Rev.\ C} \textbf{94}, 015503 (2016).

\bibitem{furmanski_pn}
  A.~P.~Furmanski and J.~T.~Sobczyk,
  ``Neutrino energy reconstruction from one-muon and one-proton events,''
  \textit{Phys.\ Rev.\ C} \textbf{95}, 065501 (2017).

\bibitem{rs_coherent}
  D.~Rein and L.~M.~Sehgal,
  ``Coherent $\pi^0$ production in neutrino reactions,''
  \textit{Nucl.\ Phys.\ B} \textbf{223}, 29--44 (1983).

\bibitem{berger_sehgal_coherent}
  C.~Berger and L.~M.~Sehgal,
  ``Partially conserved axial vector current and coherent pion production by
  low energy neutrinos,''
  \textit{Phys.\ Rev.\ D} \textbf{79}, 053003 (2009).

\bibitem{genieccpi}
  MicroBooNE Collaboration, P.~Abratenko \textit{et al.},
  ``First Measurement of Energy-Dependent Inclusive Muon-Neutrino
  Charged-Current Cross Sections on Argon with the MicroBooNE Detector,''
  \textit{Phys.\ Rev.\ Lett.} \textbf{128}, 151801 (2022), arXiv:2110.14023 [hep-ex].

\bibitem{uboone_nim}
  MicroBooNE Collaboration, R.~Acciarri \textit{et al.},
  ``Design and Construction of the MicroBooNE Detector,''
  \textit{JINST} \textbf{12}, P02017 (2017).

\bibitem{uboone_crt}
  MicroBooNE Collaboration, C.~Adams \textit{et al.},
  ``Design and construction of the MicroBooNE Cosmic Ray Tagger system,''
  \textit{JINST} \textbf{14}, P04004 (2019).

\bibitem{rs1981}
  D.~Rein and L.~M.~Sehgal,
  ``Neutrino Excitation of Baryon Resonances and Single Pion Production,''
  \textit{Ann.\ Phys.} \textbf{133}, 79--153 (1981).

\bibitem{kuzmin2004}
  K.~S.~Kuzmin, V.~V.~Lyubushkin, and V.~A.~Naumov,
  ``Lepton polarization in neutrino nucleon interactions,''
  \textit{Mod.\ Phys.\ Lett.\ A} \textbf{19}, 2815--2829 (2004).

\bibitem{graczyk2008}
  K.~M.~Graczyk and J.~T.~Sobczyk,
  ``Form factors in the quark resonance model,''
  \textit{Phys.\ Rev.\ D} \textbf{77}, 053001 (2008)
  [Erratum: \textit{Phys.\ Rev.\ D} \textbf{79}, 079903 (2009)].

\bibitem{berger2007}
  C.~Berger and L.~Sehgal,
  ``Lepton mass effects in single pion production by neutrinos,''
  \textit{Phys.\ Rev.\ D} \textbf{76}, 113004 (2007).

\bibitem{nowak2009}
  J.~A.~Nowak (MiniBooNE Collaboration),
  ``Four-momentum transfer discrepancy in the charged current $\pi^+$ production
  in the MiniBooNE: Data vs.\ theory,''
  \textit{AIP Conf.\ Proc.} \textbf{1189}, 243--248 (2009).

\bibitem{neut}
  Y.~Hayato,
  ``A neutrino interaction simulation program library NEUT,''
  \textit{Acta Phys.\ Polon.\ B} \textbf{40}, 2477--2489 (2009).

\bibitem{anl_radecky1982}
  G.~M.~Radecky \textit{et al.}\ (ANL Collaboration),
  ``Study of single pion production by weak charged currents in low-energy
  $\nu d$ interactions,''
  \textit{Phys.\ Rev.\ D} \textbf{25}, 1161--1173 (1982).

\bibitem{bnl_kitagaki1986}
  T.~Kitagaki \textit{et al.}\ (BNL Collaboration),
  ``Charged-current exclusive pion production in neutrino--deuterium
  interactions,''
  \textit{Phys.\ Rev.\ D} \textbf{34}, 2554--2565 (1986).

\bibitem{k2k_rodriguez2008}
  A.~Rodriguez \textit{et al.}\ (K2K Collaboration),
  ``Measurement of single charged pion production in the charged-current
  interactions of neutrinos in a 1.3~GeV wide band beam,''
  \textit{Phys.\ Rev.\ D} \textbf{78}, 032003 (2008).

\bibitem{miniboone_ccpi}
  A.~A.~Aguilar-Arevalo \textit{et al.}\ (MiniBooNE Collaboration),
  ``Measurement of $\nu_\mu$-induced charged-current single positive pion
  production cross sections on mineral oil at $E_\nu \in 0.5$--$2.0$~GeV,''
  \textit{Phys.\ Rev.\ D} \textbf{83}, 052007 (2011).

\bibitem{minerva_tki}
  X.-G.~Lu \textit{et al.}\ (MINERvA Collaboration),
  ``Measurement of final-state correlations in neutrino muon-proton mesonless
  production on hydrocarbon at $\langle E_\nu\rangle=3$~GeV,''
  Phys.\ Rev.\ Lett.\ \textbf{121}, 022504 (2018).

\bibitem{minerva_ccpi}
  B.~Eberly \textit{et al.}\ (MINERvA Collaboration),
  ``Charged pion production in $\nu_\mu$ interactions on hydrocarbon at
  $\langle E_\nu \rangle = 4.0$~GeV,''
  \textit{Phys.\ Rev.\ D} \textbf{92}, 092008 (2015).

\bibitem{minerva_fe}
  C.~L.~McGivern \textit{et al.}\ (MINERvA Collaboration),
  ``Cross sections for $\nu_\mu$ and $\bar{\nu}_\mu$ induced pion production on
  hydrocarbon in the few-GeV region using MINERvA,''
  \textit{Phys.\ Rev.\ D} \textbf{94}, 052005 (2016).

\bibitem{minerva_lownu}
  MINERvA Collaboration, J.~Wolcott \textit{et al.},
  ``Evidence for Quasielastic Charge Current Antineutrino Scattering Off
  Protons in Hydrocarbon at $\langle E_{\bar\nu} \rangle = 3.5$~GeV,''
  \textit{Phys.\ Rev.\ Lett.} \textbf{116}, 081802 (2016).

\bibitem{t2k_ccpip}
  K.~Abe \textit{et al.}\ (T2K Collaboration),
  ``Measurement of the muon neutrino charged-current single $\pi^+$ production
  on hydrocarbon using the T2K off-axis near detector ND280,''
  \textit{Phys.\ Rev.\ D} \textbf{101}, 012007 (2020), arXiv:1909.03936 [hep-ex].

\bibitem{argoneut_ccpi}
  R.~Acciarri \textit{et al.}\ (ArgoNEUT Collaboration),
  ``First measurement of the cross section for $\nu_\mu$ and $\bar{\nu}_\mu$
  induced single charged pion production on argon using ArgoNEUT,''
  \textit{Phys.\ Rev.\ D} \textbf{98}, 052002 (2018), arXiv:1804.10294 [hep-ex].

\bibitem{uboone_nue_ccpi}
  MicroBooNE Collaboration, P.~Abratenko \textit{et al.},
  ``First measurement of $\nu_e$ and $\bar{\nu}_e$ charged-current single
  charged-pion production differential cross sections on argon using the
  MicroBooNE detector,''
  \textit{Phys.\ Rev.\ Lett.} \textbf{135}, 061802 (2025),
  arXiv:2503.23384 [hep-ex].

\bibitem{uboone_bnb2026_ccpi}
  MicroBooNE Collaboration, P.~Abratenko \textit{et al.},
  ``Measurement of single charged pion production in charged-current
  $\nu_\mu$-Ar interactions with the MicroBooNE detector,''
  \textit{Phys.\ Rev.\ D} \textbf{113}, 032007 (2026),
  arXiv:2509.03628 [hep-ex].


\bibitem{companion_prd} \uboone{} Collaboration, ``Measurement of $\nu_\mu$ and $\bar\nu_\mu$
  charged-current single charged-pion production differential cross sections on argon in the NuMI
  beam,'' companion article, submitted to Phys.\ Rev.\ D (joint submission).
\bibitem{t2k_tki} K.~Abe \emph{et al.} (T2K Collaboration), ``Characterization of nuclear effects
  in muon-neutrino scattering on hydrocarbon with a measurement of final-state kinematics and
  correlations in charged-current pionless interactions at T2K,'' Phys.\ Rev.\ D \textbf{98},
  032003 (2018).
\bibitem{companion_prl} \uboone{} Collaboration, ``First measurement of transverse kinematic
  imbalance and hadronic-system kinematics in charged-current single-pion production on argon,''
  companion Letter, submitted to Phys.\ Rev.\ Lett.\ (joint submission).
\end{thebibliography}


\end{document}
